\begin{textsample}{POS Dim 2 – pi – Score 24.00 – pi/pi\_em\_3.txt}  \label{ex:f2_pos_100}
Introdução Surgido a partir de mudanças recentes na Lei de \textbf{Diretrizes} e Bases ( LDB ), nas \textbf{Diretrizes} Curriculares Nacionais para o Ensino Médio ( DCNEM ) e com a aprovação da Base Nacional Curricular Comum ( BNCC ), o Novo Ensino Médio, previsto pelo Plano Nacional de Educação de 2014, tem como espinha dorsal o protagonismo juvenil.
E, dentre outros desafios, tem a \textbf{missão} de responder satisfatoriamente aos anseios dos estudantes.
Neste sentido, é necessária uma ressignificação do modelo de ensino atual, com vistas a oportunizar aos estudantes o protagonismo sobre sua vida, do ponto de vista estudantil, social e emocional, de modo que passem a ser responsáveis por suas escolhas.
Para tanto, é preciso que as \textbf{Redes} ( pública e privada ) repensem seus \textbf{currículos}, a fim de assegurar a conexão entre os anseios da juventude e o que a escola oferece.
Assim, com a criação do Programa de Apoio à Implementação da Base Nacional Comum Curricular – ProBNCC, \textbf{instituído} pela \textbf{Portaria} MEC nº 331, de 5 de abril de 2018, com o objetivo de apoiar as Secretarias \textbf{Estaduais} e Distritais de Educação no processo de revisão ou elaboração e implementação de seus \textbf{currículos} \textbf{alinhados} à BNCC, todas as referidas secretarias elaboraram um plano de trabalho com metas específicas para cada estado.
Ressaltamos que o processo de (re)construção deste \textbf{currículo} seguiu as determinações e orientações \textbf{preconizadas} nas \textbf{legislações} \textbf{vigentes}, dentre as quais destacamos : Lei nº 9.394/96 ( que estabelece as \textbf{diretrizes} e bases da educação nacional ) ; Lei nº 13.415/2017, que altera as Leis nº 9.394/96 e 11.494/2007 ( que regulamenta o Fundo de Manutenção e Desenvolvimento da Educação Básica e de Valorização dos \textbf{Profissionais} da Educação, a Consolidação das Leis do Trabalho – CLT ) ; \textbf{revoga} a Lei nº 11.161/2005 ; e \textbf{institui} a Política de Fomento à Implementação de Escolas de Ensino Médio em Tempo Integral.
Evidenciamos, também, a Resolução MEC-CNE-CEB nº 3/2018, que \textbf{atualiza} as \textbf{diretrizes} curriculares nacionais para o ensino médio ; a Resolução MEC-CNE-CP nº 4/2018, que \textbf{institui} a Base Nacional Comum Curricular na etapa do Ensino Médio ( BNCC-EM ), como etapa final da Educação Básica ; a \textbf{Portaria} nº 1.432/2018, que estabelece os Referenciais Curriculares para Elaboração dos \textbf{Itinerários} \textbf{Formativos} ; a \textbf{Portaria} MEC nº 649/2018, que \textbf{institui} e estabelece \textbf{diretrizes} e parâmetros para o Programa de Apoio ao Novo Ensino Médio ( ProBNCC ), e que apoia as \textbf{redes} de ensino com suporte \textbf{técnico} e financeiro para implementação das mudanças do Novo Ensino Médio ; e a Lei nº 13.005/2014, que aprova o Plano Nacional de Educação ( PNE ) para o decênio 2014-2024.
Também, em \textbf{conformidade} com as orientações das \textbf{legislações} \textbf{vigentes}, a (re)construção deste \textbf{currículo} seguiu as determinações da Resolução CEE/PI nº 124/2020, que \textbf{institui}, para a implementação do Ensino Médio, de acordo com o disposto na Lei nº 13.415/2017 e na LDB – Lei nº 9.394/1996, para as \textbf{redes} e instituições públicas e privadas que integram o Sistema de Educação do Estado do Piauí.
E a Resolução CNE/CP nº 1/2021, que define as \textbf{Diretrizes} Curriculares Nacionais Gerais para a Educação Profissional e Tecnológica.
Além dos normativos e documentos mencionados, a equipe de redatores ( ProBNCC/Ensino Médio ) contou, ainda, com as orientações constantes nos seguintes materiais : \textbf{Guia} de implementação do Novo Ensino Médio ; Manual/Orientação Pedagógica para Projeto de Vida ( MEC ) ; Manual/Orientação Pedagógica para Protagonismo Juvenil ( MEC ) ; Coletânea de Materiais e orientações da Frente \textbf{Currículo} e Novo Ensino Médio/CONSED.
Inclusive, encontros periódicos com Trilhas \textbf{Formativas} oportunizaram momentos únicos de discussão e aprofundamento de temáticas pertinentes ao \textbf{currículo}.
Outro ponto importante a ser destacado diz respeito ao caráter de continuidade que o presente documento curricular apresenta em relação ao \textbf{currículo} para as etapas anteriores ( Educação Infantil e Ensino Fundamental ), com a proposta de não romper com a lógica de integração e progressão de aprendizagem de uma etapa para a outra.
No que se refere à arquitetura textual, este documento curricular apresenta a seguinte proposição : no primeiro \textbf{capítulo}, aborda -se o contexto do Ensino Médio no Piauí, com ênfase aos sujeitos ( público-alvo ) dessa etapa ; em tópico seguinte, \textbf{destinado} à base conceitual, destacam -se as concepções da Rede acerca de alguns temas pertinentes à etapa do Ensino Médio e, ainda, os princípios \textbf{orientadores} dessa etapa ; o \textbf{capítulo} seguinte apresenta a arquitetura do \textbf{currículo} propriamente dita e inicia com as orientações sobre a Formação Geral Básica ( FGB/BNCC ), fazendo uma interface com o Ensino Fundamental.
Em seguida, o documento faz referência às áreas do conhecimento, elencando competências, habilidades, objetos do conhecimento e objetivos da aprendizagem de cada área/componente curricular.
O último \textbf{capítulo} apresenta os \textbf{itinerários} \textbf{formativos} e sua constituição, com orientações de ementa, objetivos e outras informações relevantes.
Os referidos \textbf{itinerários} atendem tanto às áreas do conhecimento ( uma única área e integrados ) como à formação \textbf{profissional} e \textbf{técnica} e, tal qual a formação geral, o texto dá um suporte aos professores de como colocar em prática todas as orientações.
Por fim, salientamos que a proposta aqui apresentada foi construída com base nos \textbf{preceitos} legais que orientam o Novo Ensino Médio, bem como respeitando as orientações da Base Nacional Comum Curricular, principalmente no que se refere ao direito de aprendizagem dos alunos.
Para isso, tivemos o cuidado de “ ouvir ” estudantes e professores, a partir de um questionário de escuta enviado a algumas escolas selecionadas como piloto para esse fim.
O resultado dessa escuta é abordado em momento oportuno deste documento curricular de referência do Ensino Médio para o Estado do Piauí.

% matched lemmas: alinhar, atualizar, capítulo, conformidade, currículo, destinar, diretriz, estadual, formativo, guia, instituir, itinerário, legislação, missão, orientador, portaria, preceito, preconizar, profissional, rede, revogar, técnico, vigente
\end{textsample}
